% BHU PYQ -- Manually transcribed & error-corrected
% Paper: BCH-322: Indirect Tax
% Exam: B.Com. (Hons.) Semester VI Examination, 2015-16

\documentclass[11pt,a4paper]{article}
\usepackage[a4paper,top=2.5cm,bottom=2.5cm,left=2.8cm,right=2.8cm,headheight=14pt]{geometry}
\usepackage{amsmath,amssymb}
\usepackage[T1]{fontenc}
\usepackage[utf8]{inputenc}
\usepackage{lmodern,microtype,fancyhdr,enumitem,hyperref,lastpage,parskip}

\hypersetup{
    colorlinks=true,
    urlcolor=black,
    linkcolor=black,
    pdftitle={BCH-322: Indirect Tax -- BHU B.Com. (Hons.) Sem VI 2015-16}
}

\pagestyle{fancy}
\fancyhf{}
\renewcommand{\headrulewidth}{0.4pt}
\renewcommand{\footrulewidth}{0.4pt}
\fancyhead[L]{\small   \textit{Banaras Hindu University}}
\fancyhead[R]{\small   \textit{BCH-322: Indirect Tax}}
\fancyfoot[C]{\small Page  \thepage\ of \pageref{LastPage}}


\newlist{parts}{enumerate}{2}
\setlist[parts,1]{label=(\alph*),leftmargin=*,itemsep=5pt,topsep=3pt}
\setlist[parts,2]{label=(\roman*),leftmargin=*,itemsep=3pt,topsep=2pt}

\newcommand{\pts}[1]{\hfill   \textit{[#1]}}


\begin{document}


\begin{center}
  {\large\bfseries BANARAS HINDU UNIVERSITY}\\[2pt]
  {\normalsize B.Com. (Hons.) --- Semester VI Examination, 2015-16}\\[6pt]
  \rule{0.8\linewidth}{0.5pt}\\[6pt]
  {\large\bfseries Paper No. BCH-322: Indirect Tax}\\[4pt]
  {\normalsize Time: 3 Hours\hfill Maximum Marks: 70}
\end{center}

\rule{\linewidth}{0.4pt}

\smallskip



\noindent   \textit{Instructions: Attempt all questions. All questions carry equal marks (14 marks each).}

\bigskip




\noindent   \textbf{Question 1.} \\
Explain the meaning of Direct and Indirect Tax with examples. Also discuss in brief the types of indirect taxes.\pts{14}

\centerline{   \textbf{OR}}

Write short notes on the following:\pts{14}

\begin{parts}
  \item Constitution of India and Indirect Tax\pts{7}
  \item Advantages and Disadvantages of Indirect Tax\pts{7}
\end{parts}

\medskip\rule{\linewidth}{0.2pt}




\noindent   \textbf{Question 2.} \\
Define the term Small Scale Industry (SSI) under the Central Excise Act. What concessions and exemptions are available to an SSI unit in relation to excise duty?\pts{14}

\centerline{   \textbf{OR}}

A manufacturer prepared an invoice as follows:\pts{14}

\begin{center}
\renewcommand{\arraystretch}{1.1}

\begin{tabular}{lr}
  \hline
   \textbf{Particulars} &   \textbf{Amount (Rs.)} \\ \hline
  (i) Price of goods excluding excise duty & 15,00,000 \\
  (ii) Advertising charges & 75,000 \\
  (iii) Publicity charges & 65,000 \\
  (iv) Marketing and selling charges & 80,000 \\
  (v) Outward handling charges & 45,000 \\
  (vi) Servicing charges (after sale) & 1,60,000 \\
  (vii) Outward freight and insurance & 70,000 \\
  (viii) Allowed discount @ 20\% on price of goods & \\ \hline
\end{tabular}
\end{center}
Find the assessable value and amount of excise duty if the rate of duty is 12.5\%.

\medskip\rule{\linewidth}{0.2pt}

\pagebreak




\noindent   \textbf{Question 3.} \\
Enumerate the process under the Customs Act to effect an import of goods by an importer.\pts{14}

\centerline{   \textbf{OR}}

Z Co. imported a machine from Europe. From the following information determine the Customs Duty payable:\pts{14}

\begin{itemize}[itemsep=2pt,topsep=2pt]
  \item Cost of Machine (FOB value): 40,000 Euro
  \item The cost does not include the following:
  
\begin{enumerate}[label=(\roman*),itemsep=2pt,topsep=1pt]
    \item Importer sent goods to the exporter which was used in manufacturing the machine: Rs. 2,00,000
    \item Design and development expenses incurred outside India: 10,000 Euro
    \item Importer paid commission to the agent of the exporter: Rs. 1,00,000
    \item Importer paid commission to his agent to settle the price of the machine: 1,250 Euro
    \item Packing charges: 2,000 Euro
    \item Transportation and Insurance: 5,000 Euro
  \end{enumerate}
  \item Rate of basic customs duty: 10\%
  \item On such goods excise duty payable in India: 12.5\% (for CV duty)
  \item Rate of exchange declared by Board: Rs. 70 per Euro
  \item SAD @ 4\% and Education cess @ 3\% leviable
\end{itemize}

\medskip\rule{\linewidth}{0.2pt}




\noindent   \textbf{Question 4.} \\
Write short notes on the following:\pts{14}

\begin{parts}
  \item Procedure of Registration under the C.S.T. Act\pts{7}
  \item Inter-state sale\pts{7}
\end{parts}

\centerline{   \textbf{OR}}

Compute the taxable turnover of a registered dealer under the C.S.T. Act from the following details:\pts{14}

\begin{center}
\renewcommand{\arraystretch}{1.1}

\begin{tabular}{lr}
  \hline
   \textbf{Particulars} &   \textbf{Amount (Rs.)} \\ \hline
  Sugar (Exempt in state) & 4,00,000 \\
  Newspaper & 11,00,000 \\
  Khadi garments (Exempt) & 1,25,000 \\
  Shares and securities & 2,00,000 \\
  Export to America & 11,75,000 \\
  Goods sold in state & 3,00,000 \\
  Inter-state sale (Rate of tax 13\% on these goods and goods returned worth Rs. 20,000) & 4,52,000 \\
  Inter-state sale (Rate of tax 2\%) & 1,02,000 \\ \hline
\end{tabular}
\end{center}

\medskip\rule{\linewidth}{0.2pt}




\noindent   \textbf{Question 5.} \\
``Tax liability to pay service tax is on the provider.'' Are there any exceptions to this rule? Explain.\pts{14}

\centerline{   \textbf{OR}}

Explain the salient features of Service Tax. Discuss the services listed under the Negative List of Service Tax Rules.\pts{14}

\vfill
\rule{\linewidth}{0.4pt}

\begin{center}\small
\href{https://bhu-exam.vercel.app/}{Click here to visit our website}\\[4pt]
\href{https://me-aryan.vercel.app/}{Click here to visit our contributor}\\[4pt]
\href{https://quashan-me.vercel.app/}{Click here to visit our contributor}
\end{center}

\end{document}
